% !TeX program = xelatex
\documentclass[tikz, margin = 5mm]{standalone}

\usepackage{xfrac}
\usepackage{ifthen}
\usepackage{intcalc}
\usepackage[MnSymbol]{mathspec}
\setmainfont[Mapping=tex-text,Ligatures={NoRequired,NoCommon,NoContextual}]{Comic Sans MS}
\setmathfont(Digits,Latin,Greek)[Numbers={Lining,Proportional}]{Comic Sans MS}
\usetikzlibrary{shapes.geometric,decorations,decorations.pathmorphing}
\usetikzlibrary{intersections}
\usetikzlibrary{patterns}
\usetikzlibrary{calc}
\usetikzlibrary{shapes.misc}
\usetikzlibrary{spy}

\pgfdeclarelayer{bg}
\pgfdeclarelayer{bbg}
\pgfsetlayers{bbg, bg, main}

\definecolor{c1}{HTML}{ac0000}
\definecolor{c2}{HTML}{00a0d5}

\newcommand{\abs}[1]{{\left|{#1}\right|}}

\tikzset{
	point/.style n args = {1}{circle, draw = black, inner sep = #1pt, fill = black},
	empty point/.style = {circle, draw = black, inner sep = 2pt, fill = white},
	cross/.style = {cross out, draw, minimum size = 2 * (#1 - \pgflinewidth), inner sep = 0pt, outer sep = 0pt},
	xkcd/.style n args = {2}{decorate, decoration = {random steps, pre length = #1mm, post length = #1mm, segment length = #2mm, amplitude = 0.2pt}},
	xkcd/.default = {4}{0.6},
	point/.default = {2},
	declare function = {f(\x) = 0.4 * \x * sin(20/(\x) r) - 4 * cos(20 / (\x) r);}
}

\begin{document}
	\begin{tikzpicture}
		\draw[xkcd, thick, ->] (-4,0) -- (4,0) node[pos=1, below]{$x$};
		\draw[xkcd, thick, ->] (0,-4) -- (0,4) node[pos=1, left]{$y$};
		\draw[xkcd, thick, c1] (-4,-3) -- (3,4) node[pos = 1, above] {\scriptsize$f_1$};
		\draw[xkcd, thick, opacity = 0.9, c1] (-4,-3.4) -- (3.4,4) node[pos = 1, above] {\scriptsize$f_2$};
		\draw[xkcd, opacity = 0.7, c1] (-4,-3.7) -- (3.7,4) node[pos = 1, above] {\scriptsize$f_3$};
		\draw[xkcd, opacity = 0.5, c1] (-4,-3.85) -- (3.85,4) node[pos = 1, above, xshift = 0.1cm] {\scriptsize$f_4$};
		\draw[xkcd, thick, c2] (-4,-4) -- (4,4) node [pos = 1, right] {$f$};
		\node[c1] (flabel) at (-3,2.5) {$f_n(x)=x+\frac{1}{n}$};
		\draw[xkcd, c1, dashed, ->] (flabel) to[bend left] (-0.4, 0.8);
	\end{tikzpicture}
\end{document}